\section{Research Objectives and Key Scientific Questions}
\label{sec:objectives}

\subsection{Primary Research Objective}
The overarching objective of this dissertation is to design, implement, and empirically validate an end-to-end, morphologically-grounded computational framework for artificial intelligence-generated text detection and automated factual verification tailored specifically for low-resource agglutinative languages, using Kazakh as the primary benchmark language. 

By unifying symbolic morphological transducers, contrastive representation learning, hybrid lexical-neural retrieval, and a novel 2D Cartesian trust formulation, this research aims to neutralize shortcut learning pathologies, overcome subword tokenization bottlenecks, and deliver an interpretable, production-ready auditing system for low-resource NLP.

\subsection{Specific Objectives}
To achieve the primary research goal, the project is structured around four interconnected specific objectives:

\begin{enumerate}[leftmargin=*, label=\textbf{Obj \arabic*:}]
    \item \textbf{Symbolic Morphological Engineering and Feature Formulation:} Design and construct an 83-rule Finite-State Transducer (FST) rule engine and stemmer for Kazakh that accurately isolates canonical root morphemes from complex inflectional and derivational suffixation chains, and formulate a 16-dimensional continuous affix diagnostic vector ($\mathbf{m}_{\text{affix}} \in \mathbb{R}^{16}$) capturing negation mismatches, temporal conflicts, evidential modalities, and syntactic case alignments.
    \item \textbf{Morpho-Contrastive AI Text Detection Pipeline:} Develop a deep contrastive neural architecture that integrates pre-trained multilingual transformer backbones with morphological projection heads. Evaluate model invariance across a newly curated Tri-Domain corpus (formal news, encyclopedic text, and conversational social media) under strict Leave-One-Generator-Out (LOGO) cross-validation against diverse LLM families (LLaMA, GPT, Mistral, Sherkala-7B).
    \item \textbf{Construction and Validation of the Kazakh-FEVER 3K Benchmark:} Curate the first human-validated, multi-domain fact verification benchmark for Kazakh comprising 3,000 claim-evidence pairs. Establish the novel \textit{Hard Not-Enough-Info (Hard NEI)} curation protocol to enforce high lexical overlap ($\ge 85\%$) while varying truth-conditional morphological markers, and construct a hybrid evidence retrieval engine combining morpheme-stemmed BM25 with dense multilingual embeddings (\textsc{mContriever}) via Reciprocal Rank Fusion (RRF).
    \item \textbf{Mathematical Formulation of the 2D Cartesian Trust Matrix and Deployment:} Formulate a unified continuous 2D trust metric $\mathcal{T}(x) \in [0, 1]$ that jointly projects factual veracity ($T_{\text{fact}} \in [-1, 1]$) and generative authenticity ($T_{\text{gen}} \in [0, 1]$) into four operational risk quadrants, and deploy a production-grade, bilingual (Kazakh and English) interactive web platform with real-time morphemic stem highlighting for human-in-the-loop auditing.
\end{enumerate}

\subsection{Core Research Questions}
This investigation is guided and rigorously validated through four central research questions:

\begin{itemize}[leftmargin=*, label=$\blacktriangleright$]
    \item \textbf{RQ1 (Morphological Representation):} \textit{To what extent does the explicit injection of symbolic morphological diagnostic features ($\mathbf{m}_{\text{affix}}$) overcome the representation dilution caused by subword tokenizer fragmentation in agglutinative neural models?}
    \item \textbf{RQ2 (Retrieval Synergy):} \textit{Can hybrid retrieval combining rule-based morphemic BM25 indexing with dense unsupervised representations (\textsc{mContriever}) achieve superior Recall@$k$ and Mean Reciprocal Rank (MRR) across highly inflected low-resource queries compared to purely dense or purely sparse baselines?}
    \item \textbf{RQ3 (Shortcut Neutralization):} \textit{Does the Hard NEI curation protocol successfully eliminate superficial lexical overlap shortcuts in natural language inference classifiers, and does explicit morphological conditioning prevent false entailment predictions when truth-conditional allomorphs (e.g., verbal negation infixes) are manipulated?}
    \item \textbf{RQ4 (Trust Space Calibration):} \textit{Does the joint 2D Cartesian Trust Matrix $(T_{\text{fact}}, T_{\text{gen}})$ provide well-calibrated, separable operational decision boundaries across verified human ground truth, human misinformation, AI synthesis, and AI disinformation compared to standard independent thresholding?}
\end{itemize}
